\section{Two stable standard factors of comparable natural degrees}
\label{sec:two-standard-product-genus}

\paragraph{Review status and provenance.}
D-TD50, reserved question Q2ff, is a new actual-source argument by
\texttt{/root/sixfolds}.  Root suggested the almost-Galois profile
method in the Neftin--Zieve preprint and supplied the decisive
quadratic-factorization obstruction for the last almost-dihedral
case.  The unequal-degree bootstrap, Fourier-loss bounds, explicit
constants and remaining gcd estimates are reconstructed below.
G34 supplied the useful pure-kernel product-transitivity observation
following D-TD48.  The complete author proof, source-interface,
checker and first-output readings pass.  Independent whole-proof
review and root acceptance also pass, with exact scopes recorded
separately.  No preprint theorem is assumed.

\begin{theorem}[D-TD50]
Let \(C\to\mathbf P^1\) be an actual connected finite
\(J\)-Galois cover, and let its faithful absolutely irreducible
primitive complex coefficient representation have an actual
factorization
\[
 V=(\operatorname{Std}_n\circ x)\otimes
       (\operatorname{Std}_m\circ y)\otimes R,
 \qquad x(J)\in\{A_n,S_n\},\quad y(J)\in\{A_m,S_m\},
 \qquad q=\dim R\ge1.
 \tag{TPG.1}
\]
Here both natural maps and \(R\) are representations of the same
finite \(J\), and both displayed standard factor algebras are stable.
The entire remaining factor \(R\) is retained.  Suppose
\[
 m\ge2^{64},\qquad 2m-1\le n\le108m,
 \qquad e=(n-1)(m-1)q,
 \qquad \sum_i\operatorname{rank}(1-V(g_i))=2e+6
 \tag{TPG.2}
\]
for the actual original inertia tuple.  Then this configuration
does not exist.

No maximality, original character field, support floor for \(R\),
or rank-nine hypothesis is required beyond the stated actual sphere
budget.  In applications a separately proved source and decomposition
argument must supply the degree window and the two stable standard
factors.  In particular an arbitrary alternating extraction component
is not declared standard by this theorem.
\end{theorem}

\begin{proof}
\subsection*{The actual product curve and the first genus bound}

Apply accepted D-TD48 to the degree-\(n\) standard factor, keeping
\(\operatorname{Std}_m\otimes R\) as its entire complement.
The actual subgroup
\(K_0=\ker(\operatorname{Std}_m\circ y\otimes R)\)
maps injectively onto a normal subgroup containing \(A_n\), hence
acts transitively on the \(n\) letters.  It fixes every \(m\)-letter:
if \(k\in K_0\), then the tensor identity operator forces
\(\operatorname{Std}_m(y(k))\) to be scalar.  Such a scalar
standard matrix has trivial natural image for \(m\ge5\).
For example it centralizes the represented \(A_m\), so \(y(k)\)
centralizes \(A_m\) in \(S_m\), whose centralizer is trivial.

The product action on \(n m\) letters is therefore transitive:
first move the second coordinate using \(J\), and then adjust
the first coordinate using \(K_0\).  Let \(X_n,X_m,Z\) be the
actual quotient curves for the two letter stabilizers and their
intersection.  Thus \(Z\) is the connected normalization of the
full fiber product, with degrees \(m\) over \(X_n\), \(n\)
over \(X_m\), and \(nm\) over the sphere.  Write
\(g_n=g(X_n)\), \(g_m=g(X_m)\), and \(z=g(Z)\).
No independence of two abstract quotient groups is being presumed.

The accepted all-target standard estimate in HD-SS02, (SS.15),
applies separately to each factor with its entire complement.
Those complements have dimensions at least \(m-1\) and \(n-1\).
The explicit consequence \(b\ge4,N\ge17\Rightarrow g(X_N)\le2\)
therefore gives
\[
                         0\le g_n,g_m\le2.
 \tag{TPG.3}
\]
This step retains the partner cohomology; it does not discard it
or assign the original budget to either partner.

Put \(W=\operatorname{Std}_n\otimes\operatorname{Std}_m\)
with the actual maps understood, and \(d=(n-1)(m-1)\).
For the \(i\)-th original inertia let \(L_i\) be its order in
the faithful original coefficient image.  If \(f_i(j),h_i(j)\)
are the two fixed-letter counts of its \(j\)-th power, define
\[
 \chi_i(j)=(f_i(j)-1)(h_i(j)-1),\qquad
 N_i=\frac1{L_i}\sum_{j=0}^{L_i-1}\max\{-\chi_i(j),0\}.
 \tag{TPG.4}
\]
The character is real.  Fourier inversion at every target eigenvalue
and \(|\chi|=\chi+2\max(-\chi,0)\) show that its largest
eigenspace has dimension at most its one-eigenspace dimension
plus \(2N_i\).  Spectral convolution with the entire \(R\)
then gives
\[
 \operatorname{rank}(1-V(g_i))
 \ge q\bigl(\operatorname{rank}(1-W(g_i))-2N_i\bigr),
 \qquad
 z\le g_n+g_m+\sum_iN_i+\frac3q.
 \tag{TPG.5}
\]
For the second assertion, the product permutation module is
\(1\oplus\operatorname{Std}_n\oplus\operatorname{Std}_m\oplus W\).
Actual quotient pullback and Frobenius reciprocity identify the
three nonconstant multiplicities in \(H^1(C)\); hence that of
\(W\) is \(2z-2g_n-2g_m\).  Equivalently the additive
sphere Euler formula gives
\(\sum_i\operatorname{rank}(1-W(g_i))
=2d+2z-2g_n-2g_m\).
There is no invariant vector in \(W\), since the product action
is transitive.  One may also note that irreducibility of \(V\)
forces irreducibility of \(W\).  Summing the first inequality
in (TPG.5) proves its second inequality.

If both original natural permutations have a fixed letter, all
their powers do also, so \(N_i=0\).  A derangement in degree
\(N\) has index at least \(N/2\).  The natural Riemann--Hurwitz
sum is at most \(2N+2\) by (TPG.3), so at most four original
entries can be derangements in either natural action.  Thus at
most eight entries have \(N_i>0\).  At any entry
\(N_i\le c_{n,i}+c_{m,i}\le n+m\), where the \(c\)'s
are cycle counts: negative character occurs only when one fixed
count is zero, and its absolute value is at most the other fixed
count; the average fixed count over the original \(L_i\) powers
is exactly the number of cycles.  Therefore
\[
 z\le7+8(n+m)\le10(n+m),\qquad
 \mathcal R_{Z/X_n},\mathcal R_{Z/X_m}\le4096m.
 \tag{TPG.6}
\]
Indeed the two relative ramification sums are at most
\(2z+2m-2\) and \(2z+2n-2\).  The degree window gives
\(2z+2n-2\le22n+20m-2\le2396m-2<4096m\).
All Fourier averages above use the full original inertia order;
periodic fixed-count averages do not replace that order for \(V\).

\subsection*{An explicit unequal-degree local profile lemma}

Fix the constants
\[
 H=2^{12}=4096,\qquad T=2^{20},\qquad B=16T=2^{24},
 \qquad C_*=8B=2^{27},\qquad G_*=16C_*=2^{31},
 \qquad E_*=2^{54}.
 \tag{TPG.7}
\]
At one original branch value, let \(r\)'s and \(s\)'s be
the natural cycle lengths of degrees \(n,m\).  A pair of
cycles gives \(\gcd(r,s)\) product cycles.  Its relative
ramification contributions are \(r-\gcd(r,s)\) and
\(s-\gcd(r,s)\).  Each local sum is nonnegative and is at
most the corresponding total in (TPG.6).

We prove the following dichotomy at every point.  Either there
is one common integer \(k\in\{1,\ldots,6\}\) such that at
most \(B\) letters in each action lie outside its \(k\)-cycles,
or each action has at most \(T\) cycles of each length at most
six.  Call the former its mode \(k\) and the latter mode
\(\infty\).

Induct on \(k=1,\ldots,6\), assuming all previous shorter
cycle counts are at most \(T\) in both actions.  Their letter
mass is at most \(d_0=15T\).  Let \(a,b\) be the two
\(k\)-cycle counts.  Cycles longer than \(k\) satisfy
\(r-\gcd(r,k)\ge r-k\ge r/(k+1)\).  Consequently
\[
 Hm\ge\frac{b(n-ka-d_0)}{k+1},\qquad
 Hm\ge\frac{a(m-kb-d_0)}{k+1}.
 \tag{TPG.8}
\]
If \(b\le m/(2k)\), then \(d_0\le m/4\) gives
\(a\le4H(k+1)\le28H<T\).  Since
\(n-ka-d_0\ge n-21T\ge n/2\), the first inequality
then gives \(b\le2H(k+1)m/n\le14H<T\).
The induction continues.

If \(b>m/(2k)\), the first inequality gives
\(n-ka\le15T+2Hk(k+1)\le15T+84H<B\).
Thus \(a\ge n/(2k)\), and the second inequality similarly
gives \(m-kb<B\).  This is the common finite mode and ends
the induction.  All denominator estimates used here follow
already from \(m\ge2^{40}\), \(n\ge m\).
Two different finite modes cannot occur because each would
occupy more than half the letters.  If the induction never
ends early, both actions have the asserted short-cycle bounds.

\subsection*{The six possible global mode patterns}

Let \(J_n,J_m\) be the total natural indices over all mode-one
points, including zeros.  Each such point has at least \(m-B\)
fixed letters in the second action and at least \(n-B\) in
the first.  Their contributions to the relative ramification
sums give
\[
       (m-B)J_n\le Hm,\qquad (n-B)J_m\le Hm,
       \qquad J_n,J_m\le2H.
 \tag{TPG.9}
\]
There is no bound here on the number of original entries that
are identity in both natural actions, and none is needed.

In either degree \(D\in\{n,m\}\), a finite mode
\(2\le k\le6\) has index
\(D(1-1/k)+\epsilon\), with \(|\epsilon|\le B\).
An infinite mode has cycle count at most \(D/7+6T\), hence
index between \(6D/7-6T\) and \(D\).
Every non-one mode has index at least \(D/2-B\).
Five such points would exceed the natural total \(2D+2\)
when \(D>10B+4\).  There are therefore at most four.

If all are finite, the sum of their nominal weights
\(1-1/k\) must equal two.  Otherwise its distance from two
is at least \(1/60\), whereas the total error, including
(TPG.9) and the natural genus term, is at most
\(4B+2H+2\).  This is impossible at \(D\ge2^{40}\).
Enumerating the reciprocal equation with at most four
integers between two and six gives the first four rows below.

For \(l\ge1\) infinite modes, put
\(L=\sum_{\rm finite}(1-1/k)+6l/7\) and
\(U=\sum_{\rm finite}(1-1/k)+l\).
Necessarily \(L\le2\le U\): a strict failure has rational
gap at least \(1/420\), while the error is at most
\(4B+24T+2H+2\), again too small at \(D\ge2^{40}\).
Three infinite modes, or two plus any finite mode, have
\(L>2\).  One infinite mode needs exactly two finite modes;
\((2,3)\) already has \(6/7+1/2+2/3>2\), so both are
two.  Two infinite modes alone are allowed.  Thus the full
list, with all mode-one points omitted, is
\[
 (3,3,3),\quad(2,4,4),\quad(2,3,6),\quad(2,2,2,2),
 \quad(\infty,\infty),\quad(\infty,2,2).
 \tag{TPG.10}
\]
For the last two rows the nominal weights, with infinity
assigned weight one, sum to two.  Exact natural
Riemann--Hurwitz therefore bounds the sum of the infinite
cycle counts in each action by
\(4B+2H+2<C_*\).  In particular every individual such
count is at most \(C_*\).

We can now improve (TPG.6) drastically.  Mode-one points
have zero negative character.  At a finite mode \(k\ge2\),
powers divisible by \(k\) fix the bulk in both actions and
have nonnegative character.  Every other power has at most
\(B\) fixed letters in either action; its negative character
is at most \(B\).  Hence \(N_i\le B\).
At an infinite mode, the average fixed-count argument gives
\(N_i\le2C_*\).  There are at most four non-one modes, so
(TPG.5) yields
\[
                   z\le7+8C_*<G_*.
 \tag{TPG.11}
\]
This estimate still includes every original power and the
entire arbitrary factor \(R\).

\subsection*{Exact cross corrections}

At a finite mode \(k\ge2\), let \(\mathcal A,\mathcal B\)
be the multisets of exceptional cycle lengths different from
\(k\); their total letter masses are at most \(B\).
Expanding the main-main, main-exceptional and exceptional-
exceptional pairs gives the exact product cycle count
\[
 c_{nm}=
 \frac{nm}{k}
 -\frac{m}{k}\sum_{r\in\mathcal A}(r-\gcd(r,k))
 -\frac{n}{k}\sum_{s\in\mathcal B}(s-\gcd(s,k))
 +K_k,
 \quad
 K_k=\sum_{r,s}\left(
 \frac{rs-r\gcd(k,s)-s\gcd(k,r)}{k}+\gcd(r,s)\right).
 \tag{TPG.12}
\]
Each summand has absolute value at most \(4rs\), so
\(|K_k|\le4B^2\).  At a mode-one point its product index is
\[
 \operatorname{ind}(x_i\times y_i)
 =m\operatorname{ind}(x_i)+n\operatorname{ind}(y_i)-E_i,
 \quad
 E_i=\sum_{r,s}\bigl((r-1)(s-1)+\gcd(r,s)-1\bigr).
 \tag{TPG.13}
\]
Here \(0\le E_i\le2\operatorname{ind}(x_i)
\operatorname{ind}(y_i)\), since
\(\gcd(r,s)-1\le(r-1)(s-1)\).  Therefore
\(\sum_{\rm mode\ 1}E_i\le2J_nJ_m\le8H^2\).
Every sum of corrections used below has
\[
 |E|\le16B^2+8H^2<E_*.
 \tag{TPG.14}
\]
This bound permits either sign of the finite-mode corrections.

\subsection*{The four finite Euclidean patterns}

For one of the first four rows of (TPG.10), put
\[
 A_n=J_n+\sum_{\text{finite modes }k}
           \frac1k\sum_{r\in\mathcal A}(r-\gcd(r,k)),
 \qquad
 A_m=J_m+\sum_{\text{finite modes }k}
           \frac1k\sum_{s\in\mathcal B}(s-\gcd(s,k)).
 \tag{TPG.15}
\]
Both are nonnegative.  If \(A_n=0\), every mode-one
natural entry is identity and every cycle at a mode-\(k\)
point has length dividing \(k\).  The actual natural
\(A_n\)- or \(S_n\)-Galois quotient cover would then have
orbifold excess at most zero, since these nominal weights
sum to two.  Its Galois curve would have genus at most one.
This is impossible for the large natural group: the already
accepted genus-zero/genus-one finite-action argument in D-TD02
excludes such an alternating or symmetric action for degree at
least eight.  In genus zero its spherical Riemann--Hurwitz
order list makes the group dihedral or gives group order
12, 24 or 60.  In genus one the translation subgroup is
abelian with cyclic automorphism quotient of order at most six,
by Milne, Chapter~IV, Proposition~7.13 and Remark~7.18,
pp.~140--141.\cite{TD:MilneEC2021}  This is exactly the
already checked input in the complete proof of D-TD02.
Thus \(A_n>0\), and likewise \(A_m>0\).

For each individual row of (TPG.10) its denominator divides
respectively three, four, six or two.  Hence both positive
numbers are at least \(1/6\).  Summing (TPG.12)--(TPG.13)
and using the nominal weight sum two gives exactly
\[
             2z-2=mA_n+nA_m-E
                    \ge\frac{n+m}{6}-E_*.
 \tag{TPG.16}
\]

\subsection*{Two infinite points}

Let their product cycle counts be \(C_1,C_2\), and their
two cycle counts in the degree-\(n\) action be \(c_1,c_2\).
Natural and product Riemann--Hurwitz give
\[
 c_1+c_2=J_n-2g_n+2,\qquad
 2z-2=mJ_n+nJ_m-C_1-C_2-E.
 \tag{TPG.17}
\]
Both \(J_n,J_m\) are at least one.  Otherwise the actual
corresponding natural group would be generated by at most two
mutually inverse branch permutations and would be cyclic.
Since \(C_1+C_2\le m(c_1+c_2)\), (TPG.17) implies
\(2z-2\ge nJ_m+2m(g_n-1)-E_*\); the symmetric bound
also holds.  If either natural genus is positive, this is
at least \(m-E_*\).  If both genera are zero and
\(J_m\ge2\), it is at least \(2n-2m-E_*\ge2m-2-E_*\).

It remains that \(g_n=g_m=0\) and \(J_m=1\).  The
degree-\(m\) long-cycle counts sum to three, so, after relabeling
the two points, their cycle types are \([m]\) and \([a,b]\),
where \(a+b=m\) and \(1\le a\le b\).  Its only other
nonidentity natural branch permutation is a transposition.
The product-one relation identifies this with a cut of an
\(m\)-cycle at separation \(a\) or \(b\).
If \(\gcd(a,m)>1\), residues along that cycle modulo this
gcd form a nontrivial block system preserved by the cycle and
the transposition.  This contradicts natural primitivity.
Consequently \(\gcd(a,b)=1\).

Let \(c\) be the degree-\(n\) cycle count at the \([a,b]\)
point.  Then \(C_1\le m(J_n+2-c)\), and
\[
 2z-2\ge m(c-2)+n-C_2-E_*,\qquad
 C_2=\sum_{r}\bigl(\gcd(r,a)+\gcd(r,b)\bigr).
 \tag{TPG.18}
\]
The two gcds for a fixed \(r\) are coprime and their product
divides \(r\).  Hence their sum is at most \(r+1\), and
\(C_2\le n+c\).  If \(c\ge3\), (TPG.18) is at least
\(m(c-2)-c-E_*\ge m-3-E_*\).

For completeness the one- and two-cycle endpoints require
the following precise gcd estimates.  If \(b\nmid r\), then
\(\gcd(r,b)\le b/2\), so the row contributes at most
\(a+b/2\le3m/4\).  Every row contributes at most \(m\).
If \(c=2\) and either row is not divisible by \(b\), then
\(n-C_2\ge n-7m/4\ge m/4-1\).
If both are divisible, write them \(q_1b,q_2b\) and
\(q=q_1+q_2\).  Since \(n\ge2m-1>2b\), \(q\ge3\).
Coprimality gives \(C_2\le2b+q\), and therefore
\(n-C_2\ge(q-2)b-q\ge b-3\ge m/2-3\).

If \(c=1\) and \(b\nmid n\), then
\(n-m-C_2\ge m/4-1\).  Otherwise write \(n=qb\),
where again \(q\ge3\); then
\(n-m-C_2\ge(q-2)b-a-q\).
For \(q=3\), the degree window gives \(b\ge2a-1\), so
this is at least \((m-2)/3-3\ge m/4-4\).
For \(q\ge4\), it is increasing in \(q\) and is at least
\(2b-a-4\ge m/2-4\).  Thus every case of (TPG.17) gives
\[
                         2z-2\ge m/4-4-E_*.
 \tag{TPG.19}
\]

\subsection*{One infinite point and two almost-involutions}

At the two finite modes define
\[
 A_n=J_n+\frac12\sum_{\text{two points},\ r\ne2}
                          (r-\gcd(r,2)),\qquad
 t_n=\frac12\#\{\text{odd cycles at these two points}\},
 \tag{TPG.20}
\]
and similarly \(A_m,t_m\).  The \(A\)'s are nonnegative
integers.  The \(t\)'s are also integers: the number of odd
cycles of a permutation is congruent to its degree modulo two,
and there are two points.  If \(c_n,c_m\) denote the cycle
counts at the infinite point, direct natural Riemann--Hurwitz
gives
\[
 c_n=A_n+2-2g_n-t_n,\qquad
 c_m=A_m+2-2g_m-t_m,\qquad
 2z-2=mA_n+nA_m-C_\infty-E.
 \tag{TPG.21}
\]
The last identity is the exact expansion (TPG.12) at \(k=2\)
and (TPG.13) at the other points.

Both \(A_n,A_m\) are positive.  If one vanished, its
mode-one inertia would be identity and both finite-mode
permutations would be involutions or identities.  The third
permutation is their inverse product, so the entire natural
group would be generated by two involutions and be dihedral.
This is impossible for the stated natural group.

Since \(C_\infty\le m c_n\), (TPG.21) yields
\[
          2z-2\ge m(2g_n+t_n-2)+nA_m-E_*.
 \tag{TPG.22}
\]
If \(g_n\ge1\), this is at least \(n-E_*\).
If \(g_n=0\), natural primitivity forces \(t_n\ge1\).
Here is the full reason, including the higher even exceptional
cycles.  If \(t_n=0\), every ramification index in the two
finite-mode fibers is even.  The actual source curve is then
\(X_n\simeq\mathbf P^1\).  Change the coordinate on the
target so those two distinct branch values are zero and infinity.
The resulting rational function has even zero and pole divisors.
Factoring it into linear factors on \(\mathbf P^1\) shows it
equals \(c h^2\); over \(\mathbf C\), absorb a square root
of \(c\).  The natural map factors as a degree-\(n/2\)
rational map followed by a degree-two map.  Both degrees are
greater than one.  Its fibers have the resulting two nontrivial
blocks, contradicting primitivity of the natural \(A_n\) or
\(S_n\) monodromy.  This argument uses every even index,
including four, six and larger, not just transpositions.

Since \(A_m\ge1\), (TPG.22) in the remaining genus-zero
case is at least \(n-m-E_*\ge m-1-E_*\).  Thus this last
mode pattern also has a linear genus lower bound.

\subsection*{The numerical contradiction and scope}

Every mode pattern therefore implies the weaker common bound
\[
                     2z-2\ge m/6-4-E_*.
 \tag{TPG.23}
\]
At \(m\ge2^{64}\) this is strictly larger than
\(2G_*-2\), since
\(2^{64}>6\cdot2^{54}+6\cdot2^{32}+12\).
It contradicts (TPG.11).  This proves D-TD50.

The proof concerns the actual connected full product quotient,
the full original inertia powers and the entire residual factor.
No induction equality for a discarded partner, no assumed
generic monodromy and no proposed period dominance enters it.
The upper ratio bound is used only in the initial relative
ramification bound; the lower ratio bound enters the two
infinite-pattern estimates.  All constants were chosen for
transparent inequalities, not sharpness.
\end{proof}

\paragraph{Source boundary.}
Neftin--Zieve, \emph{Monodromy groups of product type},
arXiv:2403.17168v1, Proposition~4.1 and its full proof in
Section~4, motivated the local-profile method.\cite{SF:NeftinZieveProduct2024}  Its statement
has equal degrees and is not invoked here; the unequal-degree
induction and all effective constants are proved above.
The item is retained explicitly as a preprint.  The existing
accepted D-TD48, HD-SS02 and D-TD02 inputs retain their exact
published-source and review scopes.  The new quadratic
factorization and every new cycle estimate are given in full.
